\chapter{Post-nasal drip}
\index{Post-nasal drip}

\medskip
\noindent\textit{Educational information; seek professional advice for individual care.}

\section*{What it is}
Post-nasal drip is the sensation of mucus moving from the back of the nose into the throat. It is a symptom rather than a diagnosis and can occur with a cold, allergic or non-allergic rhinitis, sinus inflammation, reflux, dry air, or another nasal problem.

\section*{Symptoms}
People may notice throat clearing, a cough that is worse when lying down, a feeling of mucus or a lump in the throat, hoarseness, sore throat, bad breath, nasal blockage, or a runny nose. Similar symptoms can come from asthma, reflux, infection, or a swallowing problem.

\section*{Assessment and care}
Assessment considers duration, nasal symptoms, fever, facial pain, cough, reflux, medicines, smoking, allergies, and swallowing or breathing problems. Saline nasal rinses, fluids, humidified air, and avoiding smoke may help. Use sterile or previously boiled and cooled water for nasal irrigation and clean the device as directed. Treat the underlying allergy or infection only with appropriate advice.

\section*{When to seek medical care}
Arrange review for symptoms lasting several weeks, recurrent sinus symptoms, one-sided blockage or discharge, nosebleeds, persistent hoarseness, weight loss, or difficulty swallowing. Seek urgent help for breathing difficulty, severe headache with fever or neck stiffness, swelling around the eye, confusion, or rapidly worsening illness.

\section*{Sources}
\begin{itemize}
\item NHS allergic rhinitis information: \url{https://www.nhs.uk/conditions/allergic-rhinitis/}
\end{itemize}
